% GENERATED FILE. Edit canonical lessons, then run npm run generate:books.
% canonical-source-hash: fnv1a64:6bf46b2bb7a577d4
% canonical-lessons: RU-C410-stolknoveniye, RU-C410-stalkivatsya, RU-C410-vrezatsya, RU-C410-sbivat, RU-C410-vmyatina, RU-R410-review-403-410

\chapter{A small accident}
\label{ch:ru-a-small-accident}
\hlchaptermodality{410}

\begin{chapteropening}
\textbf{By the end of this chapter:} \emph{I can describe a small accident: a collision, colliding, crashing into something, knocking someone down and a dent.}

Complete the last lesson of chapter 410: I can describe a small accident: a collision, colliding, crashing into something, knocking someone down and a dent.
\end{chapteropening}

\section[\texorpdfstring{\ru{столкновение} (stolknovéniye)}{столкновение (stolknovéniye)}]{\ru{столкновение} (stolknovéniye) — a collision; a clash}

\label{lesson:RU-C410-stolknoveniye}



\begin{quote}
Before the new one: say the Russian for a jack, then the Russian for a car repair shop.
\end{quote}

\subsection*{You'll want to know: \ru{столкновение}}
\textbf{\ru{столкновение}} (\emph{stolknovéniye}) — \textquotedblleft{}a collision; a clash\textquotedblright{}.

\begin{quote}
\ru{На} \ru{перекрёстке} \ru{произошло} \ru{столкновение} \ru{двух} \ru{машин}.
\end{quote}

\begin{quote}
\emph{Na perekryóstke proizoshló stolknovéniye dvukh mashín.}
\end{quote}

\begin{quote}
There was a collision between two cars at the crossroads.
\end{quote}

Neuter; formal. In speech people say \emph{\ru{авария}}.

\begin{cousinweb}
\emph{\ru{столкновение}} is built on \emph{\ru{столкнуться}} 'to collide', from \emph{\ru{с}-} 'together', \emph{\ru{толкнуть}} 'to push' and \emph{-\ru{ся}}. EtymDB lists \emph{\ru{толкать}} without a deeper source.
\end{cousinweb}

\begin{grammarlens}[title={What you've built}]
One more word for telling, explaining and describing at B1.
\end{grammarlens}

\subsection*{Your turn}
\begin{itemize}
\raggedright
  \item \emph{Say it:} \emph{stolknovéniye}
  \item \emph{Say it:} \emph{stolknovéniye}, once more
  \item \emph{Say it:} a sentence of your own with \emph{stolknovéniye}
  \item \emph{Recall:} say the Russian for a jack, then the Russian for a car repair shop, then say \emph{stolknovéniye} again
\end{itemize}

\begin{tcolorbox}[breakable,colback=teal!4,colframe=teal!35!black,title={Before you move on}]
What does \ru{столкновение} mean? (\textquotedblleft{}A collision; a clash\textquotedblright{}.) Say it once more.
\end{tcolorbox}
\section[\texorpdfstring{\ru{сталкиваться} (stálkivatsya)}{сталкиваться (stálkivatsya)}]{\ru{сталкиваться} (stálkivatsya) — to collide; to come up against}

\label{lesson:RU-C410-stalkivatsya}



\begin{quote}
Before the new one: say the Russian for a car repair shop, then the Russian for a collision.
\end{quote}

\subsection*{You'll want to know: \ru{сталкиваться}}
\textbf{\ru{сталкиваться}} (\emph{stálkivatsya}) — \textquotedblleft{}to collide; to come up against\textquotedblright{}.

\begin{quote}
\ru{Зимой} \ru{на} \ru{этой} \ru{дороге} \ru{часто} \ru{сталкиваются} \ru{машины}.
\end{quote}

\begin{quote}
\emph{Zimóy na étoy doróge chásto stálkivayutsya mashíny.}
\end{quote}

\begin{quote}
In winter cars often collide on this road.
\end{quote}

Imperfective; the perfective is \emph{\ru{столкнуться}}. Takes \emph{\ru{с}} + instrumental: \emph{\ru{столкнуться} \ru{с} \ru{автобусом}}.

\begin{cousinweb}
\emph{\ru{сталкиваться}} is built from \emph{\ru{с}-} 'together', \emph{-\ru{талкивать}}, from \emph{\ru{толкать}} 'to push', and \emph{-\ru{ся}}: to push into each other. EtymDB lists \emph{\ru{толкать}} without a deeper source.
\end{cousinweb}

\begin{grammarlens}[title={What you've built}]
One more word for telling, explaining and describing at B1.
\end{grammarlens}

\subsection*{Your turn}
\begin{itemize}
\raggedright
  \item \emph{Say it:} \emph{stálkivatsya}
  \item \emph{Say it:} \emph{stálkivatsya}, once more
  \item \emph{Say it:} a sentence of your own with \emph{stálkivatsya}
  \item \emph{Recall:} say the Russian for a car repair shop, then the Russian for a collision, then say \emph{stálkivatsya} again
\end{itemize}

\begin{tcolorbox}[breakable,colback=teal!4,colframe=teal!35!black,title={Before you move on}]
What does \ru{сталкиваться} mean? (\textquotedblleft{}To collide; to come up against\textquotedblright{}.) Say it once more.
\end{tcolorbox}
\section[\texorpdfstring{\ru{врезаться} (vrézatsya)}{врезаться (vrézatsya)}]{\ru{врезаться} (vrézatsya) — to crash (into)}

\label{lesson:RU-C410-vrezatsya}



\begin{quote}
Before the new one: say the Russian for a collision, then the Russian for to collide.
\end{quote}

\subsection*{You'll want to know: \ru{врезаться}}
\textbf{\ru{врезаться}} (\emph{vrézatsya}) — \textquotedblleft{}to crash (into)\textquotedblright{}.

\begin{quote}
\ru{Водитель} \ru{не} \ru{увидел} \ru{столб} \ru{и} \ru{врезался} \ru{в} \ru{него}.
\end{quote}

\begin{quote}
\emph{Vodítel ne uvídel stolb i vrézalsya v negó.}
\end{quote}

\begin{quote}
The driver did not notice the post and crashed into it.
\end{quote}

Perfective; the imperfective \emph{\ru{врезаться}} (\emph{vrezátsya}) has the stress on the \ru{а}. Takes \emph{\ru{в}} + accusative. Future \emph{\ru{врежусь}, \ru{врежешься}}.

\begin{cousinweb}
\emph{\ru{врезаться}} is built from \emph{\ru{в}-} 'into', \emph{\ru{резать}} 'to cut' and \emph{-\ru{ся}}: to cut one's way into something. \emph{\ru{Резать}} is inherited from Proto-Slavic \emph{rězati} (EtymDB).
\end{cousinweb}

\begin{grammarlens}[title={What you've built}]
One more word for telling, explaining and describing at B1.
\end{grammarlens}

\subsection*{Your turn}
\begin{itemize}
\raggedright
  \item \emph{Say it:} \emph{vrézatsya}
  \item \emph{Say it:} \emph{vrézatsya}, once more
  \item \emph{Say it:} a sentence of your own with \emph{vrézatsya}
  \item \emph{Recall:} say the Russian for a collision, then the Russian for to collide, then say \emph{vrézatsya} again
\end{itemize}

\begin{tcolorbox}[breakable,colback=teal!4,colframe=teal!35!black,title={Before you move on}]
What does \ru{врезаться} mean? (\textquotedblleft{}To crash (into)\textquotedblright{}.) Say it once more.
\end{tcolorbox}
\section[\texorpdfstring{\ru{сбивать} (sbivát)}{сбивать (sbivát)}]{\ru{сбивать} (sbivát) — to knock down; to knock off}

\label{lesson:RU-C410-sbivat}



\begin{quote}
Before the new one: say the Russian for to collide, then the Russian for to crash.
\end{quote}

\subsection*{You'll want to know: \ru{сбивать}}
\textbf{\ru{сбивать}} (\emph{sbivát}) — \textquotedblleft{}to knock down; to knock off\textquotedblright{}.

\begin{quote}
\ru{Сильный} \ru{ветер} \ru{сбивал} \ru{людей} \ru{с} \ru{ног}.
\end{quote}

\begin{quote}
\emph{Sílnyy véter sbivál lyudéy s nog.}
\end{quote}

\begin{quote}
The strong wind was knocking people off their feet.
\end{quote}

Imperfective; perfective \emph{\ru{сбить}} (\emph{\ru{собью}, \ru{собьёшь}}). \emph{\ru{Машина} \ru{сбила} \ru{собаку}} — a car hit a dog.

\begin{cousinweb}
\emph{\ru{сбивать}} is built from \emph{\ru{с}-} 'down, off' and \emph{-\ru{бивать}}, from \emph{\ru{бить}} 'to beat'. \emph{\ru{Бить}} is inherited from Proto-Slavic \emph{biti} 'to beat' (EtymDB).
\end{cousinweb}

\begin{grammarlens}[title={What you've built}]
One more word for telling, explaining and describing at B1.
\end{grammarlens}

\subsection*{Your turn}
\begin{itemize}
\raggedright
  \item \emph{Say it:} \emph{sbivát}
  \item \emph{Say it:} \emph{sbivát}, once more
  \item \emph{Say it:} a sentence of your own with \emph{sbivát}
  \item \emph{Recall:} say the Russian for to collide, then the Russian for to crash, then say \emph{sbivát} again
\end{itemize}

\begin{tcolorbox}[breakable,colback=teal!4,colframe=teal!35!black,title={Before you move on}]
What does \ru{сбивать} mean? (\textquotedblleft{}To knock down; to knock off\textquotedblright{}.) Say it once more.
\end{tcolorbox}
\section[\texorpdfstring{\ru{вмятина} (vmyátina)}{вмятина (vmyátina)}]{\ru{вмятина} (vmyátina) — a dent}

\label{lesson:RU-C410-vmyatina}



\begin{quote}
Before the new one: say the Russian for to crash, then the Russian for to knock down.
\end{quote}

\subsection*{You'll want to know: \ru{вмятина}}
\textbf{\ru{вмятина}} (\emph{vmyátina}) — \textquotedblleft{}a dent\textquotedblright{}.

\begin{quote}
\ru{На} \ru{двери} \ru{машины} \ru{теперь} \ru{большая} \ru{вмятина}.
\end{quote}

\begin{quote}
\emph{Na dverí mashíny tepér bolsháya vmyátina.}
\end{quote}

\begin{quote}
There is now a big dent in the car door.
\end{quote}

Feminine. Built on \emph{\ru{вмять}} — to press in.

\begin{cousinweb}
\emph{\ru{вмятина}} is built on \emph{\ru{вмять}} 'to press in', from \emph{\ru{в}-} 'into' and \emph{\ru{мять}} 'to crumple, to knead' (EtymDB lists \emph{\ru{мять}} with that sense).
\end{cousinweb}

\begin{grammarlens}[title={What you've built}]
One more word for telling, explaining and describing at B1.
\end{grammarlens}

\subsection*{Your turn}
\begin{itemize}
\raggedright
  \item \emph{Say it:} \emph{vmyátina}
  \item \emph{Say it:} \emph{vmyátina}, once more
  \item \emph{Say it:} a sentence of your own with \emph{vmyátina}
  \item \emph{Recall:} say the Russian for to crash, then the Russian for to knock down, then say \emph{vmyátina} again
\end{itemize}

\begin{tcolorbox}[breakable,colback=teal!4,colframe=teal!35!black,title={Before you move on}]
What does \ru{вмятина} mean? (\textquotedblleft{}A dent\textquotedblright{}.) Say it once more.
\end{tcolorbox}
\section[(dialogue)]{Roads, stations and airports — a review}

\label{lesson:RU-R410-review-403-410}



\begin{quote}
Say the Russian for a steering wheel and a brake.
\end{quote}

\subsection*{Your turn}
Say these aloud:

\begin{itemize}
\raggedright
  \item a steering wheel, a brake, to brake, to speed up, to sound the horn
  \item a bonnet, an engine, a car battery, to start, to stall
  \item a tyre, a puncture, a pump, to pump up, to spin
  \item a boot, a bumper, a headlight, a seat, a windscreen wiper
  \item to overtake, to change lanes, to turn, to turn round, a hard shoulder
  \item speed, to exceed, an inspector, a driving licence, to fine
  \item a breakdown, a tow truck, to tow, a jack, a car repair shop
  \item a collision, to collide, to crash, to knock down, a dent
\end{itemize}

\begin{tcolorbox}[breakable,colback=teal!4,colframe=teal!35!black,title={Before you move on}]
A steering wheel? (\textbf{\ru{руль}}.) A brake? (\textbf{\ru{тормоз}}.) To brake? (\textbf{\ru{тормозить}}.) To speed up? (\textbf{\ru{разгоняться}}.) To sound the horn? (\textbf{\ru{сигналить}}.) A bonnet? (\textbf{\ru{капот}}.) An engine? (\textbf{\ru{двигатель}}.) A car battery? (\textbf{\ru{аккумулятор}}.) To start? (\textbf{\ru{заводить}}.) To stall? (\textbf{\ru{глохнуть}}.) A tyre? (\textbf{\ru{шина}}.) A puncture? (\textbf{\ru{прокол}}.) A pump? (\textbf{\ru{насос}}.) To pump up? (\textbf{\ru{накачивать}}.) To spin? (\textbf{\ru{буксовать}}.) A boot? (\textbf{\ru{багажник}}.) A bumper? (\textbf{\ru{бампер}}.) A headlight? (\textbf{\ru{фара}}.) A seat? (\textbf{\ru{сиденье}}.) A windscreen wiper? (\textbf{\ru{дворник}}.) To overtake? (\textbf{\ru{обгонять}}.) To change lanes? (\textbf{\ru{перестраиваться}}.) To turn? (\textbf{\ru{сворачивать}}.) To turn round? (\textbf{\ru{разворачиваться}}.) A hard shoulder? (\textbf{\ru{обочина}}.) Speed? (\textbf{\ru{скорость}}.) To exceed? (\textbf{\ru{превышать}}.) An inspector? (\textbf{\ru{инспектор}}.) A driving licence? (\textbf{\ru{права}}.) To fine? (\textbf{\ru{штрафовать}}.) A breakdown? (\textbf{\ru{поломка}}.) A tow truck? (\textbf{\ru{эвакуатор}}.) To tow? (\textbf{\ru{буксировать}}.) A jack? (\textbf{\ru{домкрат}}.) A car repair shop? (\textbf{\ru{автосервис}}.) A collision? (\textbf{\ru{столкновение}}.) To collide? (\textbf{\ru{сталкиваться}}.) To crash? (\textbf{\ru{врезаться}}.) To knock down? (\textbf{\ru{сбивать}}.) A dent? (\textbf{\ru{вмятина}}.)
\end{tcolorbox}
